\chapter{Courbes à terme, stockage et rendement d'opportunité}\label{ch:m3:forward-curves-storage-and-convenience-yield}

De janvier 1985 à avril 2024, le cours du contrat à terme WTI le plus proche est passé de 25.92~\$ à 86.91~\$ le
baril, soit une hausse de 3.1~\% par an. Un investisseur qui a détenu les contrats à terme sur ces mêmes années, en les renouvelant
chaque mois du contrat arrivant à échéance vers le suivant, a moins gagné sur le pétrole lui-même~: le roll
a coûté 0.7~\% par an, car le marché a passé plus de la moitié de ses jours en contango et a payé
aux acheteurs de la courbe moins que n'a monté le cours spot. La différence entre ce que fait le
prix d'une matière première et ce que rapporte une position en contrats à terme sur elle est la forme de sa \omterm{def:m3:forward-curves-storage-and-convenience-yield:curve}{courbe à terme}. Ce
chapitre explique cette forme par l'économie du stockage, mesure le \omterm{def:m3:forward-curves-storage-and-convenience-yield:storage}{rendement d'opportunité}
qu'implique le marché, montre pourquoi les contrats proches sont plus volatils que les contrats lointains, et suit
les capitaux qui suivent les \omterm{def:m3:forward-curves-storage-and-convenience-yield:index}{indices de matières premières} au fil de leur roll mensuel.

\section{La courbe à terme d'une matière première}

\begin{definition}[Courbe à terme]\label{def:m3:forward-curves-storage-and-convenience-yield:curve}
La \emph{courbe à terme}\index{courbe à terme} d'une matière première à une date est l'ensemble des prix,
à cette date, des contrats à terme ou des forwards sur la matière première pour des dates de livraison successives, considéré
comme une fonction du temps restant avant livraison.
\end{definition}

Une courbe de matière première se distingue d'une courbe de taux (le livre~2, chapitre~9) par ce qui relie ses points
entre eux~: non le seul taux d'intérêt, mais le stockage. Un baril livré dans un an est un baril détenu aujourd'hui, stocké pendant un
an dans une cuve, ou un baril produit plus tard. La courbe est en contango ou en backwardation (le livre~1,
chapitre~21), selon celle des deux façons d'obtenir le baril qui revient le moins cher.

\begin{omfigure}
\centering
\begin{tikzpicture}[font=\small]
\begin{axis}[width=11cm,height=4.8cm,
  xlabel={temps avant livraison},ylabel={prix (schématique)},
  tick label style={font=\small},label style={font=\small},
  xmin=0,xmax=12,ymin=60,ymax=100,xtick=\empty,ytick=\empty,
  legend columns=3,legend style={font=\footnotesize,at={(0.5,-0.2)},anchor=north,draw=none,fill=none,/tikz/every even column/.append style={column sep=8pt}}]
\addplot[omThm,thick,dashed,domain=0:12,samples=40] {80*exp(0.06*x/12*1.8)};
\addplot[omDef,very thick,domain=0:12,samples=40] {80*exp(0.03*x/12*1.8)};
\addplot[omProp,very thick,domain=0:12,samples=40] {80*exp(-0.12*x/12*1.8)};
\legend{portage complet~: $r+u$,contango,backwardation}
\end{axis}
\end{tikzpicture}

\omcaption{\omterm{def:m3:forward-curves-storage-and-convenience-yield:curve}{Courbes à terme}, schématiques. Le stockage plafonne la pente que peut atteindre le contango (\omterm{def:m3:forward-curves-storage-and-convenience-yield:carry}{portage complet}, c'est-à-dire le coût
de financement et de stockage), alors que rien ne plafonne la backwardation, fixée par la valeur que le marché accorde au
fait de détenir la matière première maintenant. Schématique, sans données.}\label{fig:m3:forward-curves-storage-and-convenience-yield:shapes}
\end{omfigure}

\section{La théorie du stockage et le rendement d'opportunité}

\begin{definition}[Cash-and-carry, portage complet]\label{def:m3:forward-curves-storage-and-convenience-yield:carry}
Un \emph{cash-and-carry}\index{cash-and-carry} consiste à acheter une matière première (ou un contrat à terme proche) et
à vendre un contrat à terme plus lointain en face, en finançant et en stockant la matière première dans l'intervalle. Une courbe est en
\emph{portage complet}\index{portage complet} (full carry) lorsque le prix lointain dépasse le prix proche exactement du coût
de financement et de stockage, de sorte que l'opération ne rapporte rien.
\end{definition}

Le \omterm{def:m3:forward-curves-storage-and-convenience-yield:carry}{cash-and-carry} borne le contango par le haut dès que le stockage est disponible~: $F_{t,T} \le
S_t e^{(r+u)\tau}$, avec $u$ le coût de stockage proportionnel (le plancher de stockage de
la \cref{prop:m3:crude-oil:floor} lu dans l'autre sens). Aucune opération ne borne la backwardation~: pour profiter
d'un prix lointain inférieur au prix proche, il faudrait vendre la matière première maintenant et la racheter plus tard,
ce que seul un détenteur d'inventaire peut faire, et un détenteur qui a besoin de la matière première ne le fera pas.

\begin{definition}[Théorie du stockage, rendement d'opportunité]\label{def:m3:forward-curves-storage-and-convenience-yield:storage}
La \emph{théorie du stockage}\index{théorie du stockage} (Kaldor, Working, Brennan) explique le
spread futures--spot par le coût du stockage et par un avantage à détenir de l'inventaire. Cet avantage, par
unité de valeur et de temps, est le \emph{rendement d'opportunité}\index{rendement d'opportunité} (convenience yield) $y_c$~: la valeur de
pouvoir utiliser la matière première maintenant, éviter une rupture de stock, faire tourner une usine. Avec ce rendement,
\[
F_{t,T} = S_t\,e^{(r + u - y_c)\,\tau}.
\]
\end{definition}

\begin{proposition}[Le rendement d'opportunité baisse avec l'inventaire]\label{prop:m3:forward-curves-storage-and-convenience-yield:inventory}
Entre deux contrats à terme distants de $\Delta\tau$, le \omterm{def:m3:forward-curves-storage-and-convenience-yield:storage}{rendement d'opportunité} net implicite dans le marché est
$y_c - u = r - \ln(F_2/F_1)/\Delta\tau$. Lorsque les inventaires sont abondants, $y_c$ est proche de zéro et la courbe
est proche du \omterm{def:m3:forward-curves-storage-and-convenience-yield:carry}{portage complet}~; lorsqu'ils sont rares, $y_c$ est élevé et la courbe passe en backwardation.
\end{proposition}

\begin{proof}
Prendre le logarithme de $F_2/F_1 = e^{(r + u - y_c)\Delta\tau}$. La seconde affirmation exprime le contenu économique
de la théorie~: le baril marginal en stock vaut davantage à détenir lorsqu'on en détient moins.
\end{proof}

\begin{example}[Lecture d'un spread]\label{ex:m3:forward-curves-storage-and-convenience-yield:read}
Le premier contrat WTI vaut 80.00~\$ et le second 79.40~\$, à un mois d'intervalle, avec un taux de 4~\%.
Le \omterm{def:m3:forward-curves-storage-and-convenience-yield:storage}{rendement d'opportunité} net vaut alors $0.04 - 12\ln(79.40/80.00) = 13.0\%$ par an~: le marché paie environ 13~\%
par an au-dessus du financement pour détenir du pétrole maintenant. Avec le second à 80.60~\$, il vaut $-5.0\%$~: le stockage est
payé.
\end{example}

\begin{omfigure}
\begin{tikzpicture}
\begin{axis}[width=11cm,height=4.8cm,
  ylabel={$y_c-u$, \% par an},
  tick label style={font=\small},label style={font=\small},
  xmin=1985,xmax=2024.4,ymin=-100,ymax=100,ytick={-50,0,50,100},grid=major,grid style={gray!25},
  xtick={1985,1990,1995,2000,2005,2010,2015,2020},xticklabel style={/pgf/number format/1000 sep={}},
  restrict y to domain=-100:100,unbounded coords=jump,
  table/col sep=comma]
\addplot[omDef,thin] table[x=t,y=ycu]{figdata/markets-3/10-forward-curves-storage-and-convenience-yield/convenience.csv};
\draw[gray] (axis cs:1985,0) -- (axis cs:2024.4,0);
\end{axis}
\end{tikzpicture}

\omcaption{\omterm{def:m3:forward-curves-storage-and-convenience-yield:storage}{Rendement d'opportunité} net du WTI implicite par les contrats 1 et 2 et le bon du Trésor à 3 mois,
moyennes mensuelles des valeurs quotidiennes, de janvier 1985 à mars 2024 (mois inférieurs à $-100\%$, comme avril 2020,
sortent du graphique). Le mois médian est proche de zéro~; la courbe était en backwardation pendant 45~\% des jours.
Données~: EIA~; FRED DTB3.}\label{fig:m3:forward-curves-storage-and-convenience-yield:cy}
\end{omfigure}

Lorsque le \omterm{def:m3:forward-curves-storage-and-convenience-yield:carry}{cash-and-carry} rémunère davantage que les réservoirs à terre, il rémunère aussi les navires.

\begin{definition}[Stockage flottant]\label{def:m3:forward-curves-storage-and-convenience-yield:floating}
Le \emph{stockage flottant}\index{stockage flottant} est du pétrole (ou un autre liquide) conservé dans des pétroliers à l'ancre,
en tant qu'inventaire plutôt qu'expédié, ce qui est rentable lorsque le contango sur les mois de l'affrètement dépasse
le coût de l'affrètement augmenté du financement.
\end{definition}

En plein contango de 2020, c'est ce qui s'est produit~: l'EIA, citant ClipperData, indique que le \omterm{def:m3:forward-curves-storage-and-convenience-yield:floating}{stockage flottant} mondial de brut
est passé de 65.9 millions de barils le 1\up{er} mars 2020 à 221.5 millions le 9 juillet.

\section{Rendement de roll et décomposition des rendements des contrats à terme}

Une position sur contrats à terme ne porte pas de carry propre~: son entrée ne coûte rien. Ce qu'un investisseur gagne en
détenant des contrats à terme et en les rollant provient de trois sources.

\begin{definition}[Rendement spot, rendement de roll, rendement du collatéral]\label{def:m3:forward-curves-storage-and-convenience-yield:roll}
Pour une position longue sur contrats à terme entièrement collatéralisée et rollée, le \emph{rendement spot}\index{rendement spot}
est la variation du prix du contrat le plus proche~; le \emph{rendement de roll}\index{rendement de roll} est le reste du
rendement des contrats à terme, gagné à mesure que chaque contrat détenu converge vers le prix au comptant le plus proche
(positif en backwardation, négatif en contango)~; le \emph{rendement du collatéral}\index{rendement du collatéral} est
l'intérêt perçu sur la trésorerie qui garantit la position. Leur somme constitue le rendement total.
\end{definition}

\begin{proposition}[Rendement de roll à partir de la courbe]\label{prop:m3:forward-curves-storage-and-convenience-yield:rolly}
Si la forme de la courbe reste inchangée sur une période $\Delta\tau$, une position sur un contrat situé $\Delta\tau$
plus loin que le contrat le plus proche rapporte $\ln(F_{\mathrm{near}}/F_{\mathrm{far}})/\Delta\tau$ par an au-dessus du
\omterm{def:m3:forward-curves-storage-and-convenience-yield:roll}{rendement spot}~: le \omterm{def:m3:forward-curves-storage-and-convenience-yield:roll}{rendement de roll} est égal au \omterm{def:m3:forward-curves-storage-and-convenience-yield:storage}{rendement d'opportunité} net diminué du taux,
$y_c - u - r$.
\end{proposition}

\begin{proof}
Avec une forme inchangée, le contrat détenu finit là où le contrat plus proche avait commencé~: son rendement logarithmique
est celui du spot augmenté de $\ln(F_{\mathrm{near}}/F_{\mathrm{far}})$. Il suffit d'y substituer la relation de carry.
\end{proof}

Mesurée sur les données du WTI, une position qui rolle chaque mois (comme le font les indices, dans la fenêtre décrite
ci-dessous) a gagné un \omterm{def:m3:forward-curves-storage-and-convenience-yield:roll}{rendement spot} de 3.08~\% par an de janvier 1985 à avril 2024, un \omterm{def:m3:forward-curves-storage-and-convenience-yield:roll}{rendement de roll} de $-0.71\%$
et un \omterm{def:m3:forward-curves-storage-and-convenience-yield:roll}{rendement du collatéral} de 3.18~\%~: soit un total de 5.55~\% par an
(\cref{fig:m3:forward-curves-storage-and-convenience-yield:index}).

\begin{omfigure}
\begin{tikzpicture}
\begin{axis}[width=11cm,height=5cm,ymode=log,
  ylabel={croissance de 1 (échelle logarithmique)},
  tick label style={font=\small},label style={font=\small},
  xmin=1985,xmax=2024.4,ymin=0.3,ymax=20,grid=major,grid style={gray!25},
  ytick={0.5,1,2,4,8,16},yticklabels={0.5,1,2,4,8,16},
  xtick={1985,1990,1995,2000,2005,2010,2015,2020},xticklabel style={/pgf/number format/1000 sep={}},
  legend columns=2,legend style={font=\footnotesize,at={(0.5,-0.18)},anchor=north,draw=none,fill=none,/tikz/every even column/.append style={column sep=8pt}},
  table/col sep=comma]
\addplot[omDef,thick] table[x=t,y=spot]{figdata/markets-3/10-forward-curves-storage-and-convenience-yield/index.csv};
\addplot[omThm,thick] table[x=t,y=excess]{figdata/markets-3/10-forward-curves-storage-and-convenience-yield/index.csv};
\legend{prix du contrat le plus proche,position roulée (rendement en excès)}
\end{axis}
\end{tikzpicture}

\omcaption{WTI, de janvier 1985 à avril 2024~: le prix du contrat le plus proche et une position longue roulée chaque mois sur
les 5e--9e jours de bourse, toutes deux partant de 1. L'écart est le \omterm{def:m3:forward-curves-storage-and-convenience-yield:roll}{rendement de roll} cumulé. Les
identités des contrats suivent la règle d'échéance avec un calendrier ne tenant compte que des week-ends, de sorte que quelques dates
de roll peuvent être décalées d'un jour. Données~: contrats EIA 1 à 4.}\label{fig:m3:forward-curves-storage-and-convenience-yield:index}
\end{omfigure}

\section{L'effet Samuelson}

\begin{definition}[Effet Samuelson]\label{def:m3:forward-curves-storage-and-convenience-yield:samuelson}
L'\emph{effet Samuelson}\index{effet Samuelson} est la tendance des prix des contrats à terme à être d'autant plus
volatils qu'ils sont proches de la livraison.
\end{definition}

Les chocs sur l'offre et la demande sont absorbés par l'inventaire~: un choc aujourd'hui déplace surtout le prix proche, et les
prix plus éloignés moins, car les stocks et la production ont le temps de réagir. Si l'écart du spot à son niveau de long terme
revient à la moyenne à la vitesse $\kappa$ (un processus d'Ornstein--Uhlenbeck, livre~4, chapitre~4), un contrat
dont l'échéance est à $\tau$ se déplace de $e^{-\kappa\tau}$ fois le choc du spot, et sa volatilité décroît avec $\tau$. Dans les données du WTI,
la volatilité annualisée des variations quotidiennes passe de 40.9~\% pour le premier contrat à 38.3~\%, 35.3~\% et 33.6~\%
pour le deuxième, le troisième et le quatrième
(\cref{fig:m3:forward-curves-storage-and-convenience-yield:samuelson}). Le livre~6, chapitre~16, construit sur cet effet
des modèles factoriels de la courbe entière.

\begin{omfigure}
\begin{tikzpicture}
\begin{axis}[width=8cm,height=4.4cm,ybar,bar width=18pt,
  xlabel={contrat (mois jusqu'à l'échéance, approximativement)},ylabel={volatilité annualisée, \%},
  tick label style={font=\small},label style={font=\small},
  xmin=0.4,xmax=4.6,ymin=0,ymax=48,xtick={1,2,3,4},grid=major,grid style={gray!25},
  nodes near coords,every node near coord/.append style={font=\footnotesize},
  table/col sep=comma]
\addplot[fill=omDef!60,draw=omDef] table[x=contract,y=vol]{figdata/markets-3/10-forward-curves-storage-and-convenience-yield/samuelson.csv};
\end{axis}
\end{tikzpicture}

\omcaption{Volatilité annualisée des variations logarithmiques quotidiennes des contrats WTI 1 à 4, de janvier 1985 à avril
2024, hors le jour de roll de chaque contrat et hors les jours où les prix ne sont pas strictement positifs. Données~: EIA.}\label{fig:m3:forward-curves-storage-and-convenience-yield:samuelson}
\end{omfigure}

\section{Indices de matières premières et empreinte de leur roll}

\begin{definition}[Indice de matières premières, fenêtre de roll]\label{def:m3:forward-curves-storage-and-convenience-yield:index}
Un \emph{indice de matières premières}\index{indice de matières premières} est une règle publiée pour un panier de positions sur
contrats à terme de matières premières (quels contrats, selon quelles pondérations, rollés quand), dont les investisseurs suivent
le rendement total ou en excès au moyen de swaps, de fonds et de titres de créance. Sa \emph{fenêtre de roll}\index{fenêtre de roll}
est l'ensemble des jours, fixés par la règle, où l'indice vend les contrats arrivant à échéance et achète les suivants, à parts
égales chaque jour.
\end{definition}

\begin{dated}{2026-09}{Deux fenêtres de roll d'indices}\label{dat:m3:forward-curves-storage-and-convenience-yield:windows}
Le S\&P GSCI roule sur ses 5e à 9e jours ouvrés de chaque mois, à raison de 20~\% par jour~; le
Bloomberg Commodity Index sur ses 6e à 10e jours ouvrés, également à 20~\% par jour.
\end{dated}

Une \omterm{def:m3:forward-curves-storage-and-convenience-yield:index}{fenêtre de roll} publiée signifie que chacun sait, des semaines à l'avance, que les gros capitaux indiciels vendront le contrat
proche et achèteront le suivant pendant cinq jours connus. Les traders peuvent prendre le même spread avant la fenêtre et le
dénouer pendant celle-ci, en vendant à l'indice ce qu'il doit acheter. Mou (2010) a mesuré l'impact sur les prix du roll de Goldman
de 2000 à mars 2010, a constaté que des stratégies pratiquant le front-running avec des ratios de Sharpe pouvant atteindre
4.39, et a estimé que les investisseurs indiciels avaient renoncé à 3.6~\% de rendement annuel. Plusieurs fournisseurs d'indices
proposent depuis des variantes qui rollent d'autres jours ou choisissent le contrat à détenir en fonction de la forme de la courbe.

\begin{proposition}[Le coût d'un roll prévisible]\label{prop:m3:forward-curves-storage-and-convenience-yield:rollcost}
Supposons que l'indice fasse rouler $N$ barils sur $D$ jours, par parts égales, et que la transaction de chaque jour déplace le spread
à son détriment de $k$ dollars par million de barils, ce déplacement persistant jusqu'à la fin de la fenêtre. En exécutant
après son propre impact chaque jour, l'indice paie en moyenne $k\,(N/D)\,(D + 1)/2$ par baril de plus que le spread
d'avant la fenêtre.
\end{proposition}

\begin{proof}
Le jour $j$, le spread a bougé de $j\,k\,N/D$~; la moyenne de $j$ sur $1,\dots,D$ vaut $(D + 1)/2$.
\end{proof}

Avec 50 millions de barils, cinq jours et 0.004~\$ par million de barils (à titre illustratif), le spread se déplace
de 0.04~\$ par jour et le roll coûte 0.12~\$ par baril, soit 6~millions de dollars par mois.

\section{Tutoriel~: rendement d'opportunité implicite et rendement de roll à partir de quatre contrats}\label{tut:m3:forward-curves-storage-and-convenience-yield:roll}

\begin{tutorial}
\textbf{Objectif.} À partir des quatre contrats WTI les plus proches et du taux des bons du Trésor, calculer le \omterm{def:m3:forward-curves-storage-and-convenience-yield:storage}{rendement
d'opportunité} net, les volatilités de Samuelson et le rendement d'une position roulée de type indice, décomposé en ses
parties. \textbf{État final~:}
\cref{fig:m3:forward-curves-storage-and-convenience-yield:cy,fig:m3:forward-curves-storage-and-convenience-yield:index,fig:m3:forward-curves-storage-and-convenience-yield:samuelson}.
\begin{tutsteps}
\item \textbf{Arithmétique du portage.} \omterm{def:m3:forward-curves-storage-and-convenience-yield:storage}{Rendement d'opportunité}, \omterm{def:m3:forward-curves-storage-and-convenience-yield:carry}{portage complet} (full carry), \omterm{def:m3:forward-curves-storage-and-convenience-yield:roll}{rendement de roll} et calendrier de roll.
\omcode{code/firm/commcurve/firm_commcurve.py}{36}{56}{Relations de portage et calendrier de roll de l'indice.}\label{lst:m3:forward-curves-storage-and-convenience-yield:carry}
\item \textbf{La position roulée.} Identités des contrats à partir de la règle d'échéance, pondérations à partir de la fenêtre.
\omcode{code/markets-3/10-forward-curves-storage-and-convenience-yield/python/m3_curves.py}{71}{90}{Une position WTI longue, roulée sur les 5e à 9e jours de bourse.}\label{lst:m3:forward-curves-storage-and-convenience-yield:index}
\item \textbf{Exécuter} \texttt{m3\_curves.samuelson()}, \texttt{decomposition()}, \texttt{roll\_cost()} et
\texttt{fig\_curves.py}.
\end{tutsteps}
\textbf{Pour aller plus loin.} Rouler vers le quatrième contrat au lieu du deuxième et comparer le \omterm{def:m3:forward-curves-storage-and-convenience-yield:roll}{rendement de roll}~; rouler
sur les cinq derniers jours avant l'échéance et voir ce que le prix négatif d'avril 2020 fait à la
position.
\end{tutorial}

\section{Construction~: l'objet courbe de matières premières}\label{bld:m3:forward-curves-storage-and-convenience-yield:commcurve}

\begin{build}
\bfield{Rôle} Chaque desk de matières premières de la mini-firme fixe ses prix, couvre ses positions et roule le long
d'une \omterm{def:m3:forward-curves-storage-and-convenience-yield:curve}{courbe à terme}~: l'objet courbe interpole celle-ci, lit son \omterm{def:m3:forward-curves-storage-and-convenience-yield:storage}{rendement d'opportunité}, et planifie et mesure les rolls.
\bfield{Interface} \texttt{ForwardCurve(taus, prices).price(tau)}, \texttt{.shape()}~;
\texttt{net\_convenience\_yield(f1, f2, dt, r)}~; \texttt{full\_carry\_spread}~; \texttt{roll\_yield}~;
\texttt{roll\_weights(business\_day, start, end)}~; \texttt{decompose(excess\_log, spot\_log,
collateral\_log)}.
\bfield{Règles} Les prix doivent être positifs et les échéances croissantes (un jour de prix négatif est refusé, non
interpolé)~; l'interpolation est linéaire en log-prix~; les taux sont capitalisés en continu.
\bfield{Tests d'acceptation} \texttt{code/firm/commcurve/tests/}~: interpolation et forme de la courbe~; les identités
de portage~; le calendrier de 20~\% par jour~; la décomposition.
\bfield{Pour aller plus loin} Courbes saisonnières (gaz, électricité) avec un facteur saisonnier par mois~; la courbe construite à partir de contrats à terme
aux calendriers d'échéance différents (Brent contre WTI)~; le coût de roll de l'indice mesuré sur des spreads réels.
\end{build}

\begin{omsources}
\item N.~Kaldor, ``Speculation and economic stability'', \emph{Review of Economic Studies}, 1939~;
H.~Working, ``The theory of price of storage'', \emph{American Economic Review}, 1949~;
M.~Brennan, ``The supply of storage'', \emph{American Economic Review}, 1958.
\item P.~Samuelson, ``Proof that properly anticipated prices fluctuate randomly'',
\emph{Industrial Management Review}, 1965.
\item G.~Gorton et K.~G.~Rouwenhorst, ``Facts and fantasies about commodity futures'', \emph{Financial
Analysts Journal}, 2006.
\item Y.~Mou, ``Limits to arbitrage and commodity index investment~: front-running the Goldman roll'',
2010 (SSRN 1716841).
\item S\&P Dow Jones Indices, \emph{S\&P GSCI Methodology}~; Bloomberg, \emph{BCOM} methodology.
\item EIA, contrats NYMEX WTI 1--4~; série FRED DTB3.
\item US Energy Information Administration, \emph{This Week in Petroleum}, 21 octobre 2020.
\end{omsources}

\section{Exercices}

\begin{exercise}[$\star$]\label{exo:m3:forward-curves-storage-and-convenience-yield:1}
Le spot vaut 80~\$, le taux 4~\%, le stockage 3~\% par an et le \omterm{def:m3:forward-curves-storage-and-convenience-yield:storage}{rendement d'opportunité} 5~\%. Donner le forward
à un an.
\end{exercise}

\begin{exercise}[$\star$]\label{exo:m3:forward-curves-storage-and-convenience-yield:2}
Avec un taux à 4~\% et un stockage à 3~\%, quel est le plus grand contango à un mois sur un baril à 80~\$~?
\end{exercise}

\begin{exercise}[$\star$]\label{exo:m3:forward-curves-storage-and-convenience-yield:3}
Pourquoi la backwardation peut-elle être arbitrairement forte alors que le contango ne le peut pas~?
\end{exercise}

\begin{exercise}[$\star\star$]\label{exo:m3:forward-curves-storage-and-convenience-yield:4}
Le contrat de tête vaut 70.00~\$ et le second 71.20~\$, à un mois d'intervalle~; le taux est de 5~\%. Donner le
\omterm{def:m3:forward-curves-storage-and-convenience-yield:storage}{rendement d'opportunité} net et expliquer ce qu'il signifie.
\end{exercise}

\begin{exercise}[$\star\star$]\label{exo:m3:forward-curves-storage-and-convenience-yield:5}
Une position roulée a rapporté 5.55~\% par an et le spot a monté de 3.08~\%. Avec un \omterm{def:m3:forward-curves-storage-and-convenience-yield:roll}{rendement du collatéral} de 3.18~\%,
quel a été le \omterm{def:m3:forward-curves-storage-and-convenience-yield:roll}{rendement de roll}~?
\end{exercise}

\begin{exercise}[$\star\star$]\label{exo:m3:forward-curves-storage-and-convenience-yield:6}
Dans la \cref{prop:m3:forward-curves-storage-and-convenience-yield:rollcost}, quel est le coût du roll par baril
si l'indice roule en un jour au lieu de cinq~? En dix jours~?
\end{exercise}

\begin{exercise}[$\star\star\star$]\label{exo:m3:forward-curves-storage-and-convenience-yield:7}
\emph{Programmation.} Avec \texttt{samuelson}, donner les volatilités des contrats 1 à 4 et le rapport du
quatrième au premier. Quelle vitesse de retour à la moyenne donnerait un spot de type Ornstein--Uhlenbeck pour ce rapport
sur trois mois~?
\end{exercise}

\begin{exercise}[$\star\star\star$]\label{exo:m3:forward-curves-storage-and-convenience-yield:8}
\emph{Repérer l'erreur.} «~Le pétrole a monté de 3~\% par an depuis 1985, donc un fonds suivant les contrats à terme sur pétrole a gagné
3~\% par an plus les intérêts.~»
\end{exercise}

\section{Problème~: front-running du roll}

\begin{problem}[{Problème du week-end --- un indice, sa fenêtre et les traders qui le devancent}]\label{pb:m3:forward-curves-storage-and-convenience-yield:1}
Des capitaux indiciels équivalents à 50 millions de barils de WTI roulent chaque mois du deuxième au troisième
contrat sur cinq jours connus. Chaque million de barils échangés dans une journée déplace le spread de 0.004~\$
jusqu'à la fin de la fenêtre (à titre illustratif).

\medskip\noindent\textbf{Partie I --- Le roll.}
\begin{enumerate}
\item Combien de barils l'indice roule-t-il chaque jour~?
\item De combien les transactions de chaque journée déplacent-elles le spread~?
\item Quel est le coût moyen de l'indice par baril, par rapport au spread d'avant la fenêtre~?
\item Et par mois, en dollars~?
\item Et par an~?
\end{enumerate}

\medskip\noindent\textbf{Partie II --- Le front-running.}
\begin{enumerate}[resume]
\item Quelle position un trader prend-il avant la fenêtre~?
\item Quand la ferme-t-il, et que gagne-t-il par baril~?
\item Qu'est-ce qui limite ce qu'il peut faire~?
\item Que devient le coût de l'indice lorsque davantage de traders procèdent ainsi~?
\item Pourquoi ce n'est pas de la manipulation~?
\end{enumerate}

\medskip\noindent\textbf{Partie III --- Conception.}
\begin{enumerate}[resume]
\item Que coûterait un roll réalisé en une journée~? En dix~?
\item Pourquoi les fournisseurs d'indices gardent-ils la fenêtre publique~?
\item Comment un indice qui roule vers les contrats éloignés modifie-t-il le \omterm{def:m3:forward-curves-storage-and-convenience-yield:roll}{rendement de roll}~?
\item Comment Mou a-t-il mesuré l'impact, et qu'a-t-il trouvé~?
\item Que voit l'investisseur indiciel dans ses rendements~?
\end{enumerate}

\medskip\noindent\textbf{Partie IV --- Appréciation.}
\begin{enumerate}[resume]
\item Un roll publié est-il un coût ou un service pour l'investisseur indiciel~?
\item Pourquoi le contango rend-il le roll plus coûteux pour les investisseurs indiciels, même sans front-running~?
\item Investirait-on aujourd'hui dans un \omterm{def:m3:forward-curves-storage-and-convenience-yield:index}{indice de matières premières}~? Que vérifierait-on~?
\item Énoncer le \emph{résultat nommé}~: le coût du roll prévisible par baril et par an.
\item En une phrase~: qui touche le \omterm{def:m3:forward-curves-storage-and-convenience-yield:roll}{rendement de roll}~?
\end{enumerate}
\end{problem}

\section{Questions d'entretien}

\begin{interviewq}[$\star$ \iqroles{trader}]\label{iq:m3:forward-curves-storage-and-convenience-yield:1}
Qu'est-ce que le \omterm{def:m3:forward-curves-storage-and-convenience-yield:storage}{rendement d'opportunité}, et quand est-il élevé~?
\end{interviewq}

\begin{interviewq}[$\star$ \iqroles{trader, chercheur}]\label{iq:m3:forward-curves-storage-and-convenience-yield:2}
Décomposer le rendement d'un indice de contrats à terme sur matières premières.
\end{interviewq}

\begin{interviewq}[$\star\star$ \iqroles{chercheur}]\label{iq:m3:forward-curves-storage-and-convenience-yield:3}
Pourquoi les contrats à terme proches sont-ils plus volatils que les contrats éloignés, et comment le modéliseriez-vous~?
\end{interviewq}

\begin{interviewq}[$\star\star$ \iqroles{trader}]\label{iq:m3:forward-curves-storage-and-convenience-yield:4}
La courbe est en contango marqué et les pétroliers sont bon marché. Quelle transaction envisagez-vous, et que peut-il mal se passer~?
\end{interviewq}

\begin{interviewq}[$\star\star$ \iqroles{chercheur, risque}]\label{iq:m3:forward-curves-storage-and-convenience-yield:5}
Comment testeriez-vous si la \omterm{def:m3:forward-curves-storage-and-convenience-yield:index}{fenêtre de roll} d'un indice fait bouger les prix~?
\end{interviewq}

\begin{interviewq}[$\star\star\star$ \iqroles{développeur, chercheur}]\label{iq:m3:forward-curves-storage-and-convenience-yield:6}
Construire une série continue de contrats à terme pour le backtesting. Quels choix comptent, et quels sont les pièges~?
\end{interviewq}
